\section{Incidence hexadique--octadique, opérateur d’aire et information sectorielle}
\label{rev4:ae:seccion}

La récupération géométrique du registre acquiert un contenu physique lorsqu’elle est
composée avec des observables dont la construction conserve les incidences et les
conditions d’état utilisées. Le passage d’une hexade à son octade marquée
fournit un cas précis : l’élargissement du support transforme deux
espaces d’incidence, détermine leurs multiplicités et permet de conserver
cette distinction dans un opérateur d’aire et dans le logarithme d’un état
réduit. L’exposé suivant réunit cette construction du deuxième
article de la série~\cite{hmtII} et la représentation grassmannienne
marquée de la section précédente.

Le point de départ demeure dans l’état enrichi produit par
APP--TRIT--TPK. Les deux opérations d’APP conservent leurs divisions
complètes, TRIT détermine régime et orientation, et la composition
\(U_t=\operatorname{Upd}_t\circ\operatorname{Tra}_t\circ
\operatorname{Sel}_t\) transporte et actualise l’information de position,
de retenue et de mémoire. Le prolongement compatible et l’holonomie nonadique
conservent conjointement les réalisations arithmétique et incidentielle de la
structure discrète du continu. Le registre \(K\), les coordonnées
régionales \(\pi_{\rm HMT},\varphi_{\rm HMT},e_{\rm HMT}\), la
coordonnée \(\alpha_{\rm HMT}\) et la configuration exceptionnelle proviennent
de cette construction commune. On utilise ici ses résultats et les marques
qui permettent de les prolonger vers les observables de frontière.

\subsection{Élévation d'incidence et différence de multiplicités}

Soit \(\mathcal H\) le système d'hexades de Witt sur les douze positions du registre. Fixons la coordinatisation d'incidence \(\ell:[12]\to\mathscr D\), où \(\mathscr D\) est une dodécade du code binaire de Golay étendu \(G_{24}\), de paramètres \([24,12,8]_2\). La coordinatisation transporte les hexades de \(\mathcal H\) vers le plan résiduel sur \(\mathscr D\). Cette donnée de réalisation conserve la configuration marquée déterminée dans le développement exceptionnel. Les supports de poids huit de \(G_{24}\) constituent les octades du système \(S(5,8,24)\).

\begin{proposition}[Relèvement unique d’une hexade résiduelle]
\label{rev4:ae:elevacion}
Les ensembles \(O\cap\mathscr D\), avec \(O\) une octade et
\(|O\cap\mathscr D|=6\), forment un système \(S(5,6,12)\).
Chaque hexade de ce système est l’intersection de \(\mathscr D\) avec
une unique octade. En particulier, l’hexade marquée \(H\subset\mathscr D\)
possède un relèvement déterminé \(O_H\) et une inclusion \(H\subset O_H\).
\end{proposition}
\begin{proof}
Les poids du code sont \(0,8,12,16,24\). Si
\(j=|O\cap\mathscr D|\), la somme binaire des vecteurs de support
\(O\) et \(\mathscr D\) a pour poids \(20-2j\) ; par conséquent,
\(j\in\{2,4,6\}\). Tout quintuplet \(T\subset\mathscr D\)
appartient à une unique octade. Comme son intersection avec \(\mathscr D\)
contient \(T\), cette intersection a six éléments. On obtient ainsi
une unique hexade résiduelle sur chaque quintuplet. Si deux octades relevaient
la même hexade, toutes deux contiendraient n’importe lequel de ses quintuplets et
coïncideraient par la propriété de Steiner.
\end{proof}

Dans la suite, \(\binom H2\) désigne l’ensemble des sous-ensembles à deux
éléments de \(H\). Nous conservons ce second facteur en élargissant le premier :
\begin{equation}
 \mathcal F_6=H\times\binom H2,
 \qquad
 \mathcal F_8=O_H\times\binom H2,
 \qquad
 \jmath(h,P)=(h,P).
 \label{rev4:ae:incidencias}
\end{equation}
Un élément de ces espaces est une incidence point--paire. Le point distingué parcourt d'abord l'hexade, puis son octade, tandis que la paire continue d'appartenir à l'hexade d'origine. Cette spécification fixe l'opération responsable de l'accroissement.

\begin{proposition}[Élargissement exact de l’espace des incidences]
\label{rev4:ae:incremento}
L’application \(\jmath\) est injective et son complément est
\((O_H\setminus H)\times\binom H2\). En conséquence,
\begin{equation}
 |\mathcal F_6|=6\binom62=90,
 \qquad |\mathcal F_8|=8\binom62=120,
 \qquad |\mathcal F_8\setminus\jmath(\mathcal F_6)|=30.
 \label{rev4:ae:cardinales}
\end{equation}
\end{proposition}
\begin{proof}
L’inclusion \(H\subset O_H\) conserve les deux coordonnées et prouve
l’injectivité. La partition disjointe \(O_H=H\sqcup(O_H\setminus H)\)
induit la partition du produit cartésien. Comme
\(|O_H\setminus H|=2\) et \(\binom62=15\), le complément a
\(2\cdot15=30\) éléments.
\end{proof}

Les deux multiplicités sont obtenues avant toute distribution
probabiliste sur la frontière. La différence normalisée
\((90-120)/(24\binom62)=-1/12\) conserve le signe de la comparaison
orientée. Son interprétation dépend de l’espace des incidences dans lequel
l’opération a été effectuée.

\subsection{Transport collectif et paramètre de Barbero--Immirzi}

La représentation de frontière retient les étiquettes des deux secteurs. Pour expliciter son support, écrivons \(x_j=r_{2j-1}+3r_{2j}\), avec \(r_k\in\{0,1,2\}\), en trois coordonnées de \(\mathbb Z/9\mathbb Z\). Le regroupement positionnel
\begin{equation}
 \beta_{729}(x_1,x_2,x_3)
 =\bigl(r_1+3r_2+9r_3,\ r_4+3r_5+9r_6\bigr)
 \label{rev4:ae:reagrupamiento}
\end{equation}
est une bijection vers \((\mathbb Z/27\mathbb Z)^2\). L’unicité du
développement ternaire prouve que son inverse extrait les six chiffres et
restitue les trois couples initiaux. L’ordre positionnel est préservé
avec son information de retenue ; les additions respectives restent
définies dans leurs propres groupes.

Dans le bloc \(B_\partial=\{0,\ldots,8\}^2\) du tore d'ordre vingt-sept, on distingue les liaisons orientées
\begin{align}
 E_{90}&=\{((-1,j),(0,j)),\ ((8,j),(9,j)):0\le j\le8\},\\
 E_{120}&=\{((i,-1),(i,0)),\ ((i,8),(i,9)):0\le i\le8\}.
 \label{rev4:ae:enlaces}
\end{align}
Les coordonnées sont interprétées modulo vingt-sept. Les deux ensembles
sont disjoints et contiennent dix-huit liens chacun : deux côtés de
neuf liens par direction. Les étiquettes \(90\) et \(120\) proviennent
des espaces d’incidence ; les cardinaux des ensembles de
liens sont \(18\) et \(18\).

Soient \(u_m\) les vecteurs uniformes normalisés de \(\ell^2(\mathcal F_6)\oplus\ell^2(\mathcal F_8)\), avec \(m=90,120\), et soient \(v_m\) les vecteurs uniformes normalisés correspondants de \(\ell^2(E_{90})\oplus\ell^2(E_{120})\). Par exemple, \(u_{90}=90^{-1/2}\sum_{f\in\mathcal F_6}(\delta_f,0)\) et \(v_{90}=18^{-1/2}\sum_{e\in E_{90}}(\delta_e,0)\). Chaque paire est orthonormale.

\begin{proposition}[Conservation de la structure spectrale collective]
\label{rev4:ae:colectivo}
L’application linéaire \(\mathcal J u_m=v_m\) est une isométrie
surjective entre les deux plans collectifs. Si
\[
 D_{\rm inc}=90|u_{90}\rangle\langle u_{90}|
             +120|u_{120}\rangle\langle u_{120}|,
 \qquad
 D_\partial=90|v_{90}\rangle\langle v_{90}|
             +120|v_{120}\rangle\langle v_{120}|,
\]
alors \(\mathcal J D_{\rm inc}=D_\partial\mathcal J\). Les projecteurs sectoriels et leurs combinaisons orientées sont transportés au moyen de la même application.
\end{proposition}
\begin{proof}
L’application envoie une base orthonormale sur une autre et est déterminée
par ces images. L’égalité des opérateurs vaut sur \(u_m\),
puisque les deux membres ont pour valeur \(m v_m\). Les projecteurs sont
\((120I-D)/30\) et \((D-90I)/30\) dans chaque plan, de sorte que
leurs combinaisons sont conservées par composition linéaire.
\end{proof}

Le transport précédent agit sur les modes uniformes. Les fluctuations
intrinsèques de chaque espace d’incidences demeurent dans leurs compléments
orthogonaux. La pondération aréale utilise précisément les deux étiquettes
que cette isométrie conserve.

Le paramètre de Barbero--Immirzi est déterminé à partir de ces données et de l'orientation angulaire préalablement construite. Soient \(A\) et \(C^*\) les coordonnées angulaires en degrés, dans le relèvement fixé dans la section précédente, et posons
\[
 x=\frac{\pi_{\rm HMT}A}{180},\qquad
 y=\frac{\pi_{\rm HMT}C^*}{180},\qquad
 q_\pm=e^{-x\mp y},\qquad
 S_m(q)=\frac{q^m}{1-q^{3m}}.
\]
On conserve la chambre \(x>|y|\), de sorte que \(0<q_\pm<1\).
La chaîne orientée d’un tour assigne douze secteurs à
\(\mathcal F_6\) et une couture de retour d’orientation contraire
à \(\mathcal F_8\). Son évaluation est \(12S_{90}-S_{120}\).
Le coefficient de \((1-q)^{-1}\) est déterminé par
\[
 \lim_{q\uparrow1}(1-q)(12S_{90}(q)-S_{120}(q))
 =\frac{12}{270}-\frac1{360}=\frac1{24},
\]
car \(1-q^{3m}=(1-q)\sum_{j=0}^{3m-1}q^j\).
L’évaluation conjuguée de la chaîne, jointe à la composante angulaire,
définit le paramètre structurel
\begin{equation}
 \gamma_{\rm HMT}
 =\frac{\pi_{\rm HMT}AC^*}{180}
 +12\bigl[S_{90}(q_-)-S_{90}(q_+)\bigr]
 -\bigl[S_{120}(q_-)-S_{120}(q_+)\bigr].
 \label{rev4:ae:gamma}
\end{equation}
L'origine des coefficients \(12:-1\) réside dans la chaîne orientée. La factorisation précédente détermine son coefficient de pôle. L'inversion \(C^*\mapsto-C^*\) échange les canaux et change le signe de \(\gamma_{\rm HMT}\). Nous conservons ci-dessous la branche \(C^*>0\), \(\gamma_{\rm HMT}>0\), avec la normalisation surfacique fixée par cette même orientation.

\subsection{Retour longitudinal et détermination de l’échelle aréale}
\label{rev5:longitud}

L'unité dimensionnelle de l'aire admet une construction explicite à partir du même état enrichi. Nous écrirons \(L_\alpha\) pour la longueur physique obtenue ; l'entier \(L\) qui désigne la longueur d'un registre dans la section~\ref{sec:compat-lectores-formas} conserve son sens combinatoire. La distinction sépare le cardinal d'une représentation, l'étendue d'une arête et le rayon de sa réalisation.

La composition utilise les coordonnées régionales et le registre déjà
engendrés par APP--TRIT--TPK. APP conserve les deux feuillets avec leurs restes
et quotients ; TRIT oriente le régime local ; TPK prolonge l’état,
en conservant retenue, frontière, route et mémoire. La chaîne
\(w_6\to w_{12}\to w_{18}\to w_{24}\to w_{30}\to R_{36}\to G_9\)
et l’ordre \(U_t\to\Gamma_9\to K_{\rm ph}\) maintiennent le
retour de phase avec accroissement de mémoire. La construction conjointe
du continu conserve ses aspects spectral, solénoïdal, de jauge,
cohomologique et déterminantal. La longueur qui suit est une lecture
dimensionnelle ultérieure de cette structure et de ses sorties corrélées
\(\pi_{\rm HMT},\varphi_{\rm HMT},e_{\rm HMT},\alpha_{\rm HMT}\).

Fixons le registre marqué
\[
 \Omega=\{(j,k,q,u,v):j,k\in\mathbb Z/12\mathbb Z,\quad
 k-j\in\{0,3,6,9\},\quad1\le q\le8,\quad
 u<v,\quad u,v\in\{1,2,4,5,7,8\}\}.
\]
La soustraction est calculée dans ce groupe cyclique. Son cardinal est
\(N_\Omega=12\cdot4\cdot8\binom62=5760\).
Le facteur \(120\) provient de l’espace octadique des incidences
point--paire déjà construit ; les douze positions et les quatre
positions de chaque orbite conservent leurs marques respectives. Cette
résolution élargit le porteur du retour de Hadamard. Dans les espaces
hermitiens de départ et d’arrivée, de dimension \(N_\Omega\), soit
\(B\) l’inverseur transporté, avec
\(B^*B=BB^*=I/4\). L’identité locale
\(H_4^*H_4=4I\) produit cette normalisation pour l’inverseur ;
son extension orthogonale et le transport des marques la conservent.

Dans l'espace d'arrivée, soient \(u_\Omega=N_\Omega^{-1/2}\mathbf1\), \(P=u_\Omega u_\Omega^*\), \(C=(I-P)B\) et \(M=PB\). Le transport complet \(C+M=B\) est conservé. Les projecteurs atomiques \(P_i=e_ie_i^*\), déterminés par l'inscription, définissent l'espérance diagonale \(\Delta_\Omega(T)=\sum_iP_iTP_i\).

\begin{proposition}[Réponse enregistrée et complément conservé]
\label{rev5:retorno}
Parmi les applications linéaires vers l’algèbre diagonale qui fixent
les \(P_i\) et sont bimodulaires par rapport à cette algèbre,
\(\Delta_\Omega\) est unique. En outre,
\begin{equation}
 \Delta_\Omega(CC^*)=r_\Omega I,
 \qquad r_\Omega=\frac{N_\Omega-1}{4N_\Omega}
 =\frac{5759}{23040},
 \qquad \Delta_\Omega(MM^*)=\frac{I}{4N_\Omega}.
 \label{rev5:retorno-escalar}
\end{equation}
Les deux réponses ont pour somme \(I/4\).
\end{proposition}
\begin{proof}
Pour l’unité matricielle \(E_{ij}\), la bimodularité impose
\(E(E_{ij})=P_iE(E_{ij})P_j\). Comme l’image est diagonale,
ce terme vaut zéro si \(i\ne j\) ; si \(i=j\), fixer
\(P_i\) détermine son image. Les unités matricielles constituent
une base et déterminent l’application. D’autre part,
\(CC^*=(I-P)/4\), \(MM^*=P/4\), et chaque coefficient diagonal
de \(P\) est \(1/N_\Omega\). Appliquer l’espérance démontre
les trois identités.
\end{proof}

L'inscription isométrique \(We_i=e_i\otimes e_i\) réalise \(W^*(T\otimes I)W=\Delta_\Omega(T)\). L'inscription complète conserve l'information du registre ; l'espérance extrait sa réponse diagonale. Le complément \(M\) demeure dans le transport et possède la réponse spécifique de \eqref{rev5:retorno-escalar}. Cette séparation identifie l'observable qui intervient dans la longueur et la composante complémentaire conservée.

La forme normale de Smith de \(H_4\) est
\(\operatorname{diag}(1,2,2,4)\) ; son conoyau a pour ordre
seize. Le produit du scalaire engendré
\(\alpha=\alpha_{\rm HMT}\), indexé par les classes de ce
conoyau, définit le lecteur multiplicatif fini
\[
 n_H:=\prod_{\bar v\in\operatorname{coker}_{\mathbb Z}(H_4)}
 \alpha=\alpha^{16}.
\]
Soit \(L_*>0\) une section de référence
radiale de la droite dimensionnelle de longueur. La composition longitudinale
agit linéairement sur une cochaîne orientée \(\beta_*\) :
\begin{equation}
 \mathscr T_L\beta_*
 =\alpha^{16}(\Delta_\Omega(CC^*)\otimes I)\beta_*
 =q_\alpha\beta_*,\qquad
 q_\alpha=\alpha^{16}r_\Omega,\qquad
 L_\alpha=q_\alpha L_*.
 \label{rev5:longitud-generada}
\end{equation}
La formule fixe une composition concrète de lecteurs. La linéarité conserve le type de longueur et son orientation ; le coefficient provient du retour enregistré et de la norme locale. Par exemple, la réponse diagonale d'une source \(a e_i\) est \(a r_\Omega\), tandis que la norme \(\|C^*ae_i\|=|a|\sqrt{r_\Omega}\) constitue une autre lecture du même transport.

Si \(\beta_*(\bar a)=-U(a)\beta_*(a)\), multiplier par
\(q_\alpha\) conserve l’inversion orientée. Si une arête est
subdivisée en neuf tronçons compatibles,
\(\beta_*(a)=\sum_jU_j^{-1}\beta_*(a_j)\), la même
multiplication conserve cette identité et la concaténation de mémoire.
Dans un élargissement auxiliaire du registre, on transporte
\(P\), \(CC^*\) et \(\Delta_\Omega\) par tensorisation
avec l’identité. Ainsi, \(r_\Omega\) demeure ; le cardinal d’un
maillage raffiné et le cardinal du registre marqué ont des fonctions
distinctes.

\subsection{Réciprocité radiale et reconnaissance de la longueur de Planck}
\label{rev5:radial}

La lecture circulaire du calendrier de cent huit avancées conserve
\(108\ell_0=2\pi_{\rm HMT}L_\alpha\), avec
\(\ell_0=ct_0\). Par conséquent,
\begin{equation}
 t_0=\frac{\pi_{\rm HMT}L_\alpha}{54c},\qquad
 \omega=\frac{c}{L_\alpha},\qquad
 Q_{\rm clk}=\hbar_{\rm ret}\omega
 =\frac{\hbar_{\rm ret}c}{L_\alpha}.
 \label{rev5:reloj}
\end{equation}
L’action retournée est reçue de son équation HMT complète dans
\eqref{eq:mass-k-action} ; la vitesse est reçue des deux canaux
constitutifs de la section~\ref{sec:compat-lectores-formas}.
La relation circulaire convertit l’avancée relevée en rayon. Si
\(c=\widehat c\,u_C\), \(t_0=u_T\) et
\(u_L=u_Cu_T\), il en résulte
\(L_*/u_L=54\widehat c/(\pi_{\rm HMT}q_\alpha)\).
La référence radiale \(L_*\) et la base d’arête \(u_L\)
sont reliées avec leurs types dimensionnels explicites.

Dans une fibre finie, soit \(X=X^*>0\) un caractère positif inversible, \(U=2B\) et \(Y=UXU^*\). Le transport longitudinal est \(D_L=L_\alpha U\). Sa réponse radiale positive est définie au moyen de la métrique d'arrivée, \(\|R_Xv\|^2=\|XD_L^*v\|^2\) pour tout \(v\). L'énergie et la longueur conjuguée conservent le même caractère :
\[
 H_X=Q_{\rm clk}Y,\qquad M_X=H_X/c^2,\qquad
 \Lambda_X=L_\alpha Y^{-1}.
\]

\begin{theorem}[Coefficient radial commun et échelle aréale]
\label{rev5:rigidez-radial}
Les conditions précédentes déterminent une unique réponse positive
\(R_X=L_\alpha Y\). On a
\begin{equation}
 R_X\Lambda_X=L_\alpha^2I,\qquad
 H_X\Lambda_X=\hbar_{\rm ret}cI,\qquad
 R_X=\frac{L_\alpha^2}{\hbar_{\rm ret}c}H_X.
\end{equation}
Dans la réalisation physique qui identifie \(R_X\) au rayon
gravitationnel réduit \(R_X=(G/c^2)M_X\), le coefficient est
unique et commun à tous les caractères positifs admissibles :
\begin{equation}
 G=\frac{c^3L_\alpha^2}{\hbar_{\rm ret}},\qquad
 \frac{\hbar_{\rm ret}G}{c^3}=L_\alpha^2.
 \label{rev5:G-area}
\end{equation}
\end{theorem}
\begin{proof}
La polarisation de l'égalité des normes détermine \(R_X^2=D_LX^2D_L^*=L_\alpha^2UX^2U^*\). L'opérateur \(L_\alpha UXU^*\) est positif et son carré est celui indiqué. L'unicité de la racine positive prouve la première affirmation. Les deux identités de produits découlent d'une substitution. La convention physique donne \(L_\alpha Y=(GQ_{\rm clk}/c^4)Y\). L'inversibilité de \(Y\) permet de le simplifier, et l'expression de \(Q_{\rm clk}\) fournit \eqref{rev5:G-area}. Si un autre coefficient conserve les mêmes données, leur différence multipliée par \(M_X\) vaut zéro et est donc nulle.
\end{proof}

Le domaine de la preuve est la fibre positive finie ; il admet aussi
des opérateurs bornés uniformément positifs avec les mêmes identités.
La lecture du rayon gravitationnel réduit constitue la convention
physique déclarée de cette réalisation. La reconnaissance conventionnelle
ultérieure \(\ell_P^2=\hbar_{\rm ret}G/c^3\) identifie
\(\ell_P=L_\alpha\). La longueur s’obtient d’abord au moyen de
\eqref{rev5:longitud-generada} ; cette identification conserve son
constructeur. La formule circulaire équivalente est
\[
 G=\left(\frac{54}{\pi_{\rm HMT}}\right)^2
 \frac{c^5t_0^2}{\hbar_{\rm ret}}.
\]
Le facteur circulaire est déjà inclus lorsque l’on utilise \(L_\alpha\).
La longueur et la section d’action déterminent également
\[
 t_P=\frac{L_\alpha}{c},\qquad
 m_P=\frac{\hbar_{\rm ret}}{cL_\alpha},\qquad
 E_P=\frac{\hbar_{\rm ret}c}{L_\alpha}.
\]
Ces identités découlent de la substitution de \eqref{rev5:G-area} dans les
expressions conventionnelles des trois échelles ; la positivité
détermine leurs racines. Pour un mode de caractère \(y>0\), les
lectures sont \(m(y)=m_Py\), \(r_g(y)=L_\alpha y\) et
\(\bar\lambda(y)=L_\alpha/y\). Leur produit longitudinal
vaut \(r_g(y)\bar\lambda(y)=L_\alpha^2\). L’égalité des
deux longueurs équivaut à \(y=y^{-1}\) et, par positivité,
caractérise exactement \(y=1\). Le calendrier conserve
\(t_0=(\pi_{\rm HMT}/54)t_P\) et
\(108t_0=2\pi_{\rm HMT}t_P\).

Un rayon d’horizon \(R_h=L_\alpha\zeta_h\) conserve son
facteur propre \(\zeta_h=R_h/L_\alpha\). Son aire spatiale
appartient à l’horizon déclaré ; l’échelle élémentaire de l’opérateur
sectoriel suivant est \(L_\alpha^2\). Dans la spécialisation
de Schwarzschild pour la même masse, le rayon est \(2r_g(y)\)
et, par conséquent, \(\zeta_h=2y\) ; un horizon d’une autre réalisation
conserve sa condition géométrique de formation.


\subsection{Occupations de frontière et reconstruction sectorielle}

Chaque lien de \(E_{90}\sqcup E_{120}\) détermine un facteur
\(\mathbb C^3\), de base \(|0\rangle,|1\rangle,|2\rangle\)
et d’opérateur nombre \(N_e=\operatorname{diag}(0,1,2)\). L’espace des
états de frontière est
\(\mathscr H_A=(\mathbb C^3)^{\otimes18}\otimes
(\mathbb C^3)^{\otimes18}\). Nous définissons \(N_m\) comme la somme
des opérateurs nombre des dix-huit liens du secteur \(m\).
Les deux opérateurs commutent et chacun a un spectre entier entre
zéro et trente-six.

L’échelle angulaire provient de l’horloge et du caractère normalisé du
triplet conjugué :
\begin{equation}
 \theta_{\rm clk}=\frac{C^*}{6}\frac{\pi_{\rm HMT}}{180},
 \quad
 \chi_{10}=\frac13\Tr\operatorname{diag}
 (1,e^{i\pi_{\rm HMT}/18},e^{-i\pi_{\rm HMT}/18})
 =\frac{1+2\cos(\pi_{\rm HMT}/18)}3,
 \quad a_0=\chi_{10}\theta_{\rm clk}.
 \label{rev4:ae:escala}
\end{equation}
L'identité du caractère s'obtient en additionnant les phases conjuguées. Dans la branche considérée, \(a_0>0\). L'unité d'aire est \(L_\alpha^2=\ell_P^2\), construite et reconnue dans \eqref{rev5:longitud-generada}--\eqref{rev5:G-area}. Écrivons \(a_\partial=\gamma_{\rm HMT}a_0>0\). L'opérateur d'aire est
\begin{equation}
 \widehat{\mathcal A}=L_\alpha^2a_\partial N,
 \qquad N=N_{90}+N_{120},
 \qquad D=90N_{90}+120N_{120}.
 \label{rev4:ae:area}
\end{equation}
Son spectre est \(\{nL_\alpha^2a_\partial:0\le n\le72\}\).
La multiplicité de \(n\) est le coefficient de \(z^n\) dans
\((1+z+z^2)^{36}\), car chaque facteur tensoriel apporte l’une des
trois occupations.

\begin{theorem}[Reconstruction des deux occupations sectorielles]
\label{rev4:ae:recuperacion}
La paire d’observables \((N,D)\) détermine
\begin{equation}
 N_{90}=4N-\frac{D}{30},
 \qquad N_{120}=\frac{D}{30}-3N.
 \label{rev4:ae:inversa}
\end{equation}
Dans le spectre conjoint, ses valeurs \((n,d)\) satisfont \(n\in\mathbb Z\), \(d\in30\mathbb Z\), \(90n\le d\le120n\) et \(0\le4n-d/30\le36\), \(0\le d/30-3n\le36\). Ces conditions caractérisent l'image des occupations entières.
\end{theorem}
\begin{proof}
La transformation de \((N_{90},N_{120})\) en \((N,D)\) a pour matrice
\(\left(\begin{smallmatrix}1&1\\90&120\end{smallmatrix}\right)\),
dont le déterminant est \(30\). Son inverse donne les identités
d’opérateurs. La commutativité permet de les appliquer dans chaque espace propre
commun. Les inégalités expriment la non-négativité et les capacités
maximales des secteurs ; la divisibilité rend les deux occupations entières.
Réciproquement, chaque entier entre zéro et trente-six s’exprime comme
somme de dix-huit termes de \(\{0,1,2\}\), et donc toute paire
satisfaisant ces conditions appartient au spectre conjoint.
\end{proof}

L’aire conserve la somme des occupations ; l’observable pondéré conserve
leur composition sectorielle. Ainsi, \((N_{90},N_{120})=(1,0)\) et \((0,1)\)
ont la même aire \(L_\alpha^2a_\partial\) et des poids \(90\) et \(120\).
Dans le sens complémentaire, \((4,0)\) et \((0,3)\) ont le même poids
\(D=360\), mais des aires différentes. Le couple sépare les deux situations.
À l’intérieur de chaque espace propre conjoint, les dégénérescences
microscopiques correspondant aux distributions entre liens sont conservées.

\begin{figure}[htbp]
\centering
\begin{tikzpicture}[>=Latex,every node/.style={font=\small},node distance=15mm]
\node[draw,rounded corners,align=center] (a) {Occupation sectorielle\\$(1,0)$};
\node[draw,rounded corners,align=center,right=54mm of a] (b) {Occupation sectorielle\\$(0,1)$};
\node[draw,rounded corners,align=center] (n) at ($(a)!0.5!(b)+(0,-1.8)$)
 {Aire commune\\$L_\alpha^2a_\partial$};
\draw[->] (a)--node[left] {$N=1$}(n);
\draw[->] (b)--node[right] {$N=1$}(n);
\node[above=5mm of a] {$D=90$};
\node[above=5mm of b] {$D=120$};
\end{tikzpicture}
\caption{L'égalité d'aire conserve l'occupation totale. La pondération d'incidence distingue les deux compositions ; leur lecture conjointe permet de reconstruire les occupations des deux secteurs.}
\label{rev4:ae:figura}
\end{figure}

\subsection{État réduit, purification et entropie d'intrication}

L'interprétation informationnelle requiert de fixer l'état. Pour \(\Delta\in\mathbb R\) fini, avec la normalisation du paramètre modulaire de la construction sectorielle, nous définissons
\begin{equation}
 x_m=e^{-m\Delta},\qquad z_m=1+x_m+x_m^2,
 \qquad Z=z_{90}^{18}z_{120}^{18},
 \qquad \rho_A(\Delta)=Z^{-1}e^{-\Delta D}.
 \label{rev4:ae:rho}
\end{equation}
La somme des occupations commutantes et la décomposition tensorielle prouvent
l’expression de \(Z\). Toutes les valeurs propres de \(\rho_A\) sont
strictement positives pour \(\Delta\) réel fini. Son hamiltonien
modulaire, noté \(K_A\) pour le distinguer du registre arithmétique
\(K\), est
\begin{equation}
 K_A=-\log\rho_A=(\log Z)I+\Delta D.
 \label{rev4:ae:modular}
\end{equation}
Cet opérateur caractérise la distribution spectrale de l’état réduit.
Le générateur d’évolution temporelle est défini dans le développement dynamique
avec ses propres paramètres et son propre domaine. Pour \(\Delta\ne0\), lorsque
\(\Delta\) et la normalisation sont connus, la paire
\((\widehat{\mathcal A},K_A)\) retrouve \((N,D)\) et, au moyen de
\eqref{rev4:ae:inversa}, les deux occupations. Pour \(\Delta=0\),
l’état est uniforme et \(K_A=36\log3\,I\).

Nous construisons explicitement la bipartition. Soit \(\mathscr H_{\bar A}\) une copie de \(\mathscr H_A\), avec les bases conjuguées fixées. Pour une liaison de poids \(w\), le vecteur
\[
 |\operatorname{TFD}(w)\rangle
 =\frac1{\sqrt{1+e^{-w}+e^{-2w}}}
 \sum_{n=0}^{2}e^{-wn/2}|n\rangle_A\otimes|n\rangle_{\bar A}
\]
a une norme égale à un. L'état pur
\begin{equation}
 |\Psi_\Delta\rangle
 =|\operatorname{TFD}(90\Delta)\rangle^{\otimes18}
  \otimes|\operatorname{TFD}(120\Delta)\rangle^{\otimes18}
 \label{rev4:ae:tfd}
\end{equation}
détermine la bipartition \(A\mid\bar A\).

\begin{proposition}[Réduction sectorielle et entropie de la purification]
\label{rev4:ae:entropia}
La trace partielle de \(|\Psi_\Delta\rangle\langle\Psi_\Delta|\)
sur \(\mathscr H_{\bar A}\) est \(\rho_A(\Delta)\). Étant définis
\[
 f_m=\frac{x_m+2x_m^2}{z_m},
 \qquad s(\Delta)=-\Tr(\rho_A(\Delta)\log\rho_A(\Delta)),
\]
on a
\begin{align}
 \langle N\rangle_\Delta&=18(f_{90}+f_{120}),\\
 \langle D\rangle_\Delta&=18(90f_{90}+120f_{120}),\\
 \langle\widehat{\mathcal A}\rangle_\Delta
 &=18L_\alpha^2a_\partial(f_{90}+f_{120}),\\
 s(\Delta)&=18\log z_{90}+18\log z_{120}
             +\Delta\langle D\rangle_\Delta.
 \label{rev4:ae:entropia-formula}
\end{align}
La dernière expression est l'entropie d'intrication de la purification \eqref{rev4:ae:tfd}, en unités sans dimension.
\end{proposition}
\begin{proof}
Dans la trace partielle d’un lien, l’orthogonalité de
\(|n\rangle_{\bar A}\) élimine les termes d’indices distincts et
laisse les poids \((1,x_m,x_m^2)/z_m\). La trace partielle d’un produit
tensoriel est le produit des traces partielles ; on obtient
\eqref{rev4:ae:rho}. Le premier moment de chaque lien est \(f_m\).
L’additivité des opérateurs nombre donne les trois premières identités.
La dernière résulte de la substitution de \eqref{rev4:ae:modular} dans
\(s=\Tr(\rho_AK_A)\). Les coefficients de la construction précédente
sont des coefficients de Schmidt de la bipartition déclarée, de sorte que
cette entropie coïncide avec son entropie d’intrication.
\end{proof}

L'entropie maximale de cet espace est \(36\log3\), atteinte en \(\Delta=0\). La quantité \(36\log3-s(\Delta)\) est l'entropie relative de \(\rho_A(\Delta)\) par rapport à \(I/3^{36}\). L'incidence fixe les secteurs et leurs multiplicités ; le paramètre de l'état fixe les probabilités ; la coordinatisation dimensionnelle convertit l'occupation en aire. L'identification physique de \(\mathscr H_{\bar A}\) conserve la bipartition et l'algèbre des observables qui viennent d'être spécifiées.

\subsection{Identité modulaire finie et transport de l'information géométrique}

Fixons l'état de référence \(\rho_0=\rho_A(\Delta_0)\), avec \(\Delta_0\) réel fini, et maintenons constants \(\Delta_0,a_\partial,L_\alpha\). Les aires sectorielles normalisées \(\mathcal A_m^{\rm n}=a_\partial N_m\) satisfont
\begin{equation}
 K_0=-\log\rho_0=c_0I+
 \beta_{90}\mathcal A_{90}^{\rm n}
 +\beta_{120}\mathcal A_{120}^{\rm n},
 \qquad c_0=\log Z(\Delta_0),\qquad
 \beta_m=\frac{m\Delta_0}{a_\partial}.
 \label{rev4:ae:betas}
\end{equation}
La relation \(\beta_{120}=\tfrac43\beta_{90}\) conserve le rapport des incidences. Pour \(\Delta_0\ne0\), le rapport des deux coefficients est défini et vaut \(4/3\).

\begin{theorem}[Identité finie d'aire et d'information relative]
\label{rev4:ae:ley-finita}
Pour tout état normalisé \(\sigma\) de \(\mathscr H_A\), soit \(\delta_\sigma\langle B\rangle=\Tr[(\sigma-\rho_0)B]\). Alors
\begin{equation}
 s(\sigma)-s(\rho_0)
 =\sum_{m\in\{90,120\}}\beta_m
       \delta_\sigma\langle\mathcal A_m^{\rm n}\rangle
   -D(\sigma\Vert\rho_0),
 \label{rev4:ae:ley-finita-formula}
\end{equation}
où \(D(\sigma\Vert\rho_0)=
\Tr[\sigma(\log\sigma-\log\rho_0)]\) est finie et non négative,
avec la convention \(0\log0=0\).
\end{theorem}
\begin{proof}
La définition donne \(D(\sigma\Vert\rho_0)=-s(\sigma)+\Tr(\sigma K_0)\), tandis que \(s(\rho_0)=\Tr(\rho_0K_0)\). Soustraire les identités et substituer \eqref{rev4:ae:betas} prouve \eqref{rev4:ae:ley-finita-formula} ; le terme scalaire disparaît du fait de la normalisation des traces. Pour la positivité, soient \(p_i\) et \(u_i\) les valeurs propres et les vecteurs propres de \(\sigma\), et \(q_j>0,v_j\) ceux de \(\rho_0\). Avec \(r_i=\langle u_i,\rho_0u_i\rangle\), la concavité de \(\log\) donne \(\langle u_i,\log\rho_0u_i\rangle\le\log r_i\). Ainsi, \(D(\sigma\Vert\rho_0)\ge\sum_i p_i\log(p_i/r_i)\ge0\), par convexité de \(t\log t\), puisque \(\sum_i r_i=1\). Le rang plein de \(\rho_0\) garantit la finitude.
\end{proof}

Sa version infinitésimale maintient l’état de référence fixe. Si
\(\sigma(\varepsilon)\) est différentiable, normalisé et
\(\sigma(0)=\rho_0\), alors
\[
 \left.\frac{\dd s(\sigma(\varepsilon))}{\dd\varepsilon}\right|_0
 =\sum_m\beta_m
  \left.\frac{\dd\langle\mathcal A_m^{\rm n}\rangle_{
                          \sigma(\varepsilon)}}{\dd\varepsilon}\right|_0.
\]
En effet, la dérivée de la trace de \(-\sigma\log\sigma\) est \(-\Tr[\dot\sigma(\log\rho_0+I)]\) ; la normalisation élimine \(\Tr\dot\sigma\) et la substitution de \(K_0\) donne la formule. L'identité finie ajoute à cette réponse tangente la correction exacte d'entropie relative. Les coefficients des aires physiques sont \(\beta_m/L_\alpha^2\), en conservant la même coordinatisation dimensionnelle.

La composition avec la représentation grassmannienne précise enfin quelle information est conservée. Dans \(\mathcal F(K,\xi)=([B(K)],Q,\xi)\), les coordonnées supplémentaires \(\xi\) retiennent les grandeurs régionales, l'orientation angulaire, l'inclusion marquée, l'identification des secteurs, la base des occupations, le paramètre \(\Delta\) et l'unité aréale pertinente. L'inverse de la proposition~\ref{prop:compat-inversa-marcada} restitue \(K\) ; les autres conditions sont conservées par leurs marques. D'après le théorème~\ref{thm:compat-algebra-lectores}, la composition avec cet inverse transporte \(\gamma_{\rm HMT}\), \(\widehat{\mathcal A}\), \(\rho_A\), \(K_A\) et leurs évaluations sur le domaine déclaré. Les changements de coordinatisation du même point marqué conservent ces objets, et les raffinements conservent les évaluations héritées lorsqu'ils maintiennent les mêmes données sectorielles.

Cette articulation conserve trois niveaux distinguables d’information.
La représentation positive marquée récupère le registre avec les données
retenues ; le couple aire--hamiltonien modulaire récupère les deux
occupations sous les conditions de \eqref{rev4:ae:modular} ; la
purification détermine l’intrication d’une bipartition concrète.
La perte d’information dans une projection se calcule par ses fibres :
l’aire identifie les occupations de même somme, tandis que la lecture
conjointe sépare leur composition et conserve la dégénérescence microscopique
restante. Le lien entre géométrie et information est ainsi exprimé
par des applications et des identités vérifiables sur la même construction.
\subsection{Naturalité de l’échelle et formes canoniques opératorielles}
\label{rev5:naturalidad-areal}

La longueur engendrée permet de préciser la compatibilité entre la
réalisation positive, l’opérateur d’aire et l’information sectorielle.
La représentation marquée \(\mathcal F(K,\xi)\) conserve,
parmi les données pertinentes de \(\xi\), la résolution
\(\Omega\), ses projecteurs, le retour transporté et la section
\(L_*\). Les coordonnées régionales engendrées, l’action et les canaux
constitutifs maintiennent leur généalogie antérieure. L’inverse marqué
restitue ainsi tous les arguments de
\eqref{rev5:longitud-generada} et de \eqref{rev5:G-area}.
En particulier, le lecteur géométrique déjà construit de
\(\alpha^{\mathrm{geom}}\) détermine les expressions
\[
 L_\alpha^{\mathrm{geom}}
 =L_*(\alpha^{\mathrm{geom}})^{16}\frac{5759}{23040},
 \qquad
 G^{\mathrm{geom}}
 =\frac{(c^{\mathrm{geom}})^3
 (L_\alpha^{\mathrm{geom}})^2}{\hbar_{\rm ret}^{\mathrm{geom}}}.
\]
Les lecteurs de vitesse et d’action s’obtiennent par la même
composition avec l’inverse marqué. Ces expressions conservent
les données du retour et la section radiale de référence.

\begin{corollary}[Conservation par changements de coordinatisation et raffinement hérité]
\label{rev5:conservacion}
Les changements positifs de coordinatisation conservant le point marqué et les raffinements restituant le registre initial conservent \(L_\alpha\), \(G\), \(\gamma_{\rm HMT}\) et \(\widehat{\mathcal A}=L_\alpha^2a_\partial N\) lorsqu'ils transportent également les données sectorielles et dimensionnelles déclarées. Le spectre normalisé \(\widehat{\mathcal A}/L_\alpha^2=a_\partial N\), ses multiplicités et l'état \(\rho_A(\Delta)\) demeurent identiques.
\end{corollary}
\begin{proof}
Le théorème~\ref{thm:compat-algebra-lectores} transporte chaque lecteur
par composition avec l’inverse. Les mutations du même point
reconstruisent le même registre. Dans le raffinement hérité, la somme
des séparations subdivisées reconstruit chaque coordonnée initiale.
La tensorisation de l’inscription conserve \(r_\Omega\),
et ses autres arguments restent fixés par les marques. Par
substitution, la longueur, le coefficient gravitationnel et
l’opérateur d’aire sont conservés. Les projecteurs d’occupation conservent le
polynôme de multiplicités \((1+z+z^2)^{36}\) et l’expression
normalisée de l’état.
\end{proof}

Le raffinement hérité considéré ici conserve le même espace
de trente-six liens de frontière. Une extension qui introduit
un espace physique de frontière différent requiert des opérateurs d’entrelacement
spécifiés pour \(N_{90}\) et \(N_{120}\), ainsi que le
transport des données sectorielles déclarées.

Soit maintenant \(P\) un polytope de rang un de ce manuscrit, avec une subdivision orientée \(P=\bigcup_\sigma P_\sigma\) du même support, de multiplicité un et à faces compatibles. On fixe la même fibre sectorielle \(\mathscr H_A\) sur toutes ses cellules. La forme de degré maximal à valeurs opératorielles et de dimension d'aire est définie par
\[
 \boldsymbol\Omega^{\mathcal A}_P
 =\widehat{\mathcal A}\otimes\Omega_P.
\]
Ici \(\widehat{\mathcal A}\) est constant par rapport aux
coordonnées différentielles du polytope : il provient de l’état HMT
fixé, tandis que \(\Omega_P\) décrit son support géométrique.
Le coefficient \(\widehat{\mathcal A}\) est semi-défini positif.
Sa composante dans le secteur du vide \(N=0\) est nulle. Sur un
espace propre d’occupation \(n>0\), diviser par
\(nL_\alpha^2a_\partial\) récupère la forme canonique
\(\Omega_P\) multipliée par l’identité de cet espace propre.

\begin{proposition}[Additivité aréale et résidus de frontière]
\label{rev5:forma-areal}
Dans la trivialisation sectorielle commune, on a
\begin{equation}
 \boldsymbol\Omega^{\mathcal A}_P
 =\sum_\sigma\widehat{\mathcal A}\otimes\Omega_{P_\sigma},
 \qquad
 \operatorname{Res}_F\boldsymbol\Omega^{\mathcal A}_P
 =\widehat{\mathcal A}\otimes\Omega_F.
\end{equation}
Les identités persistent sous subdivisions successives et résidus itérés. Pour un état fixé \(\rho\), prendre la trace donne \(\Tr(\rho\widehat{\mathcal A})\Omega_P\) et conserve les mêmes identités scalaires.
\end{proposition}
\begin{proof}
L’espace d’opérateurs de \(\mathscr H_A\) est de dimension
finie. Dans toute base, l’assertion est l’identité
d’additivité des formes canoniques multipliée par chaque coefficient
constant de \(\widehat{\mathcal A}\). Le résidu et la trace
sont linéaires. L’annulation des facettes intérieures conserve leurs
orientations et le même coefficient opératoriel des deux côtés.
L’itération applique ces opérations à chaque subdivision et à chaque
drapeau de faces.
\end{proof}

Si différentes cellules portent des coefficients opératoriels réguliers \(A_\sigma\), le défaut d'une paire sur sa face commune est \((A_\sigma|_F-A_\tau|_F)\otimes\Omega_F\). L'égalité des restrictions opératorielles à la frontière fournit le recollement, sous les conditions de pôles simples du théorème~\ref{thm:compat-pegado}. La somme globale conserve en outre toutes les contributions sur chaque diviseur ambiant et le contrôle des pôles extérieurs. La normalisation dimensionnelle, l'annulation des résidus et la couverture du support ont ainsi des critères séparés et compatibles.

L’identité modulaire finie conserve aussi une expression directe
en aires physiques \(\mathcal A_m=L_\alpha^2a_\partial N_m\) :
\begin{equation}
 s(\sigma)-s(\rho_0)
 =\sum_{m=90,120}\frac{m\Delta_0}{L_\alpha^2a_\partial}
   \Tr[(\sigma-\rho_0)\mathcal A_m]
  -D(\sigma\Vert\rho_0).
 \label{rev5:modular-fisica}
\end{equation}
C’est \eqref{rev4:ae:ley-finita-formula} après insertion de la
longueur construite. Le facteur angulaire \(a_0\), contenu dans
\(a_\partial=\gamma_{\rm HMT}a_0\), demeure intégral.
Si l’on compare des sections dimensionnelles avec \(L'_\alpha=sL_\alpha\)
et si l’on conserve le même état et les mêmes données adimensionnelles,
les aires sont multipliées par \(s^2\) et leurs coefficients modulaires
par \(s^{-2}\). L’entropie et l’entropie relative demeurent
égales. Cette comparaison de sections ne constitue pas une autre solution
pour les mêmes données longitudinales fixées.

Pour un horizon sphérique de rayon \(R_h=L_\alpha\zeta_h\),
l’aire géométrique vaut \(4\pi_{\rm HMT}L_\alpha^2\zeta_h^2\).
Son identification à l’espérance de l’opérateur sectoriel, si elle est
requise pour une même coupe, exige la condition explicite
\[
 a_\partial\Tr(\rho N)=4\pi_{\rm HMT}\zeta_h^2.
\]
La géométrie de l'horizon et l'état de frontière doivent satisfaire cette relation avec leurs propres constructions. Le facteur \(\zeta_h\) caractérise cet horizon ; l'échelle longitudinale du couplage reste \(L_\alpha\).

